\section{Conclusioni}

\subsection{Esito dello studio}
Dalle fonti raccolte emerge che l'app risponde a un'esigenza reale. Nel questionario (5 studenti e 2 docenti) il problema più segnalato è il tempo necessario per trascrivere a casa o in laboratorio i dati presi sul campo (6 rispondenti su 7), e tutti si trovano in zone offline a volte o spesso. Le funzioni più apprezzate sono foto e note vocali collegate alle coordinate GPS (media 4,4 su 5), la ricerca rapida delle specie senza internet (4,1) e l'uso offline (3,8); la chat di gruppo è la meno importante (2,6). L'analisi dei mock-up ha individuato criticità di usabilità risolvibili con interventi mirati, come l'assenza dell'ora di salvataggio e il pulsante SOS che copre altri elementi.
